Computational resources let a Physics Library article point to reviewed,
precomputed files that help readers explore an example.  Typical resources are
static explorers, source archives, provenance files, licenses, and CSV data.
They are meant to support the mathematical article; they are not a separate
article type.

\section{What authors should expect}

A computational resource is attached to a saved article by uploading a small
manifest file named
\begin{verbatim}
computational-resources.json
\end{verbatim}
through the normal article file attachment controls.  The filename matters.
Do not rename it.

After the article is saved, a \emph{Computational Resources} section appears on
the saved article page.  The section appears outside the rendered article body,
near the view-style controls and article metadata.  It remains visible when a
reader switches between HTML, PDF, page images, and source views.

The editor preview only renders the current LaTeX body.  It does not show the
saved article's file attachments or the surrounding article page.  Therefore it
is normal for the computational resources section to be absent from preview and
to appear only after saving the article.

\section{What computational resources are}

Computational resources are reviewed, static publication artifacts.  They may
include:
\begin{itemize}
\item an offline HTML explorer;
\item a downloadable source archive;
\item CSV result data;
\item a provenance file recording the software/runtime used to produce the
      results;
\item a license notice.
\end{itemize}

Viewing the article or moving controls in an explorer should not start a
calculation on the Physics Library server.  Published explorers should use saved
results.  This keeps the site responsive, inexpensive to host, and resistant to
automated requests that would otherwise waste computation.

\section{How to attach an approved resource}

For an already approved computational resource, the maintainer will provide a
manifest file.  For example:
\begin{verbatim}
{
  "schema_version": 1,
  "resources": ["julia-oscillator"]
}
\end{verbatim}

To attach it to an article:
\begin{enumerate}
\item Open the article editor.
\item In the attachment area, choose the file named
      \texttt{computational-resources.json}.
\item Upload it through the existing filebox controls.
\item Save or update the article normally.
\item Open the saved article page and check the \emph{Computational Resources}
      section.
\end{enumerate}

The section does not require a new LaTeX command in the article body.  Removing
the manifest from the article's attachments and saving the article removes the
section.

\section{Writing article text that uses a resource}

The article body should still explain the mathematics.  The resource should be
presented as supporting material, not as a substitute for the argument.

Good article text usually includes:
\begin{itemize}
\item the equation or model being studied;
\item the parameters used in the published computation;
\item what was held fixed and what was varied;
\item what the saved figures or data show;
\item how a reader can reproduce the calculation from the source archive;
\item any assumptions or limitations of the model.
\end{itemize}

For example, an article might say that an explorer displays saved trajectories
for a damped harmonic oscillator and that changing the displayed case does not
solve the differential equation in the browser.

\section{Licensing convention}

Physics Library article text, mathematical exposition, and figures should keep
the normal Physics Library article license, currently Creative Commons
Attribution-ShareAlike.  Computational source code should use a software
license.  For Physics Library computational bundles, the recommended default is
the GNU General Public License, version 3 or later.

A typical bundle should therefore say:
\begin{itemize}
\item source code, build scripts, tests, and browser scripts: GPLv3 or later;
\item article-style text, mathematical exposition, plots, and result data:
      Physics Library CC BY-SA unless otherwise stated.
\end{itemize}

This split keeps software under a software license while preserving the normal
share-alike treatment of Physics Library mathematical content.

\section{What authors should not do}

\begin{itemize}
\item Do not paste arbitrary external URLs into the manifest.
\item Do not upload code and expect the server to run it.
\item Do not use article preview as the final check for resources; save and
      inspect the article page.
\item Do not rename \texttt{computational-resources.json}.
\item Do not publish unexplained numbers without enough context for a reader to
      reproduce or evaluate them.
\end{itemize}

\section{Checklist before asking for review}

\begin{enumerate}
\item The article explains the model, parameters, and limitations.
\item The resource is static and does not require server-side execution.
\item The source archive contains the code needed to reproduce the published
      results.
\item The provenance file records the runtime and important file hashes.
\item The license notice distinguishes software from article/math materials.
\item The saved article shows the \emph{Computational Resources} section.
\item The explorer, source archive, provenance, license, and CSV links all open.
\end{enumerate}

When in doubt, keep the article readable without the resource.  A good
computational resource should make the article more useful, not make the article
dependent on a particular browser feature or live computation service.
